% Quarantine source-aware proof for the componentwise rigid-exception
% endpoint.  Date: 29 August 2026.
% This file uses the separately frozen exact local zero-atom lemma.

\begin{proof}[Proof of
\cref{prop:p21p3-componentwise-rigid-exceptions}]
We separate the reductions because only the last, measure-theoretic one uses
\cref{lem:p21p3-componentwise-normal-rigid-zero-atom}.

\paragraph{Nilreduction and the ambient geometry.}
An analytic seminorm kills every nilpotent.  Hence $X$ and
$X_{\mathrm{red}}$ have the same underlying topological space, the same
analytic points, and the same completed residue fields.  Affinoid charts of
the reduction are quotients of affinoid charts of $X$ by their nilradicals.
It follows directly that goodness, separatedness, strictness, and pure
dimension one pass to $X_{\mathrm{red}}$.  The immersion $\iota$ is finite
and has empty relative boundary by \cite[\S11.4.4]{CLD25}.  The
boundary-composition identity \cite[(11.4.1.3)]{CLD25} therefore shows that
$X_{\mathrm{red}}$ is boundaryless.

Both $X$ and $X_{\mathrm{red}}$ are good and topologically separated, hence
locally holomorphically separated by \cite[\S15.2.2]{CLD25} and locally
differentially separated by \cite[Corollary~15.2.4]{CLD25}.  Thus the Radon
measure and Borel-multiplication statements below apply by
\cite[Lemma~7.2.4, Proposition~7.2.8, and Example~7.2.12]{CLD25}.

We next justify the full calculus assertion, including its topology.  If
$A\to A/I$ is an affinoid quotient with $I$ nilpotent, every unit of $A/I$
lifts to a unit of $A$: if $ab=1+n$ with $n$ nilpotent, the finite geometric
series for $(1+n)^{-1}$ corrects $b$ to an inverse.  Every analytic seminorm
on $A$ kills $I$, so a lifted unit and its reduction have exactly the same
absolute value at every analytic point.  Consequently a moment map on a
reduced chart lifts, after shrinking, to a moment map on $X$ with the same
tropicalization.  Conversely every moment on $X$ restricts to one on the
reduction.  This is the nilpotent-immersion case of
\cite[Proposition~15.1.3(b) and the following form variant]{CLD25}.

The construction of ordinary CLD forms from moment charts now gives
\begin{equation}
 \label{eq:p21p3-componentwise-form-isomorphism}
 \iota^*:A_X^{p,q}\xrightarrow{\ \sim\ }
                    A_{X_{\mathrm{red}}}^{p,q}.
\end{equation}
The equality of tropicalizations makes
\eqref{eq:p21p3-componentwise-form-isomorphism} compatible with wedge
product, $J$, $\mathrm d'$, $\mathrm d''$, support, and each positivity
cone.  For the topological assertion, fix a compact test support.  A
compact-test datum is a finite collection of moment charts and Lagerberg
representatives.  Restriction and the lifted charts just constructed are
cofinal collections of such data on the two spaces.  On corresponding
charts the tropical images, polyhedral subdivisions, representatives, and
every $C^r$-seminorm are identical.  Thus
\begin{equation}
 \label{eq:p21p3-componentwise-test-homeomorphism}
 \iota^*:A^{p,q}_{X,c}\xrightarrow{\ \sim\ }
                    A^{p,q}_{X_{\mathrm{red}},c}
\end{equation}
is a homeomorphism of locally convex spaces.

Taking continuous duals in
\eqref{eq:p21p3-componentwise-test-homeomorphism} gives
\begin{equation}
 \label{eq:p21p3-componentwise-current-isomorphism}
 \iota_*:D^{a,b}(X_{\mathrm{red}})
       \xrightarrow{\ \sim\ }D^{a,b}(X),
 \qquad
 \langle\iota_*S,\alpha\rangle
       =\langle S,\iota^*\alpha\rangle.
\end{equation}
The test-form positivity cones agree, so this isomorphism preserves
positivity.  It preserves $\mathrm d'$- and $\mathrm d''$-closedness because
$\iota^*$ commutes with the two differentials.  It also identifies the weak
topologies and multiplication by smooth functions.

If $S$ is positive, its product with a locally bounded Borel function is
defined through the Radon measures attached to positive test forms, as in
\cite[Example~7.2.12]{CLD25}.  Those test forms and attached measures agree
under \eqref{eq:p21p3-componentwise-form-isomorphism}, so
$\iota_*(fS)=f\iota_*S$.  If locally $S=S_1-S_2$ with $S_1,S_2$ positive,
the same formula defines $fS$ by linearity.  It is independent of the
decomposition: an equality
$S_1+S'_2=S'_1+S_2$ of positive currents is an equality of all attached
measures, which remains true after multiplication by $f$.  This argument
does not define a Borel product for an arbitrary order-zero current.

Finally, the first stage in the CLD psh filtration is the same on $X$ and
$X_{\mathrm{red}}$, because smoothness and positivity of
$\mathrm d'\mathrm d''u$ agree.  Successor stages take decreasing limits,
the continuous filtration takes local uniform limits, and limit stages take
local unions.  The open sets and their real-valued functions are the same.
Transfinite induction therefore identifies the psh and continuous psh
sheaves.  The $G$-form assertion follows identically from the $G$-form
variant after \cite[Proposition~15.1.3]{CLD25}.  This proves
\cref{itm:p21p3-componentwise:calculus}.

\paragraph{Normalization and the locally finite rigid locus.}
Let
\[
 \bar\nu:\coprod_i\widetilde X_i\longrightarrow X_{\mathrm{red}}
\]
be the componentwise normalization.  On a good affinoid domain define
\[
 \mathcal Q:=\bar\nu_*\mathcal O_{\coprod_i\widetilde X_i}
                  /\mathcal O_{X_{\mathrm{red}}},
 \qquad
 \mathfrak c:=\operatorname{Ann}_{\mathcal O_{X_{\mathrm{red}}}}
                  (\mathcal Q).
\]
Let $Z_0$ be the conductor locus
\[
 \operatorname{Supp}(\mathcal Q)
   =V(\mathfrak c)
   =\operatorname{Supp}(\mathcal O_{X_{\mathrm{red}}}/\mathfrak c),
\]
and let $Z$ be the union of $Z_0$ with all pairwise intersections of
distinct irreducible components.  Analytic local rings on a good analytic
space are excellent by \cite[\S11.1.10]{CLD25}; normalization is finite and
its components correspond to the irreducible components by
\cite[\S11.3.1]{CLD25}.  Thus $\mathcal Q$ and the conductor are coherent.
Their support misses the generic point of every reduced component and is
zero-dimensional.  The same dimension argument applies to a proper
intersection of two components.  Irreducible components are locally finite
by \cite[\S11.2.5]{CLD25}; on a relatively compact affinoid neighborhood
only finitely many occur.  Hence $Z$ is closed and locally finite.

Strictness turns this zero-dimensional set into a rigid set.  Locally a
zero-dimensional closed analytic subspace of a strict affinoid is the
spectrum of a zero-dimensional strict affinoid algebra $B$.  The algebra is
noetherian by \cite[Theorem~4.5]{Tate71Rigid}, hence Artinian.  The strict
affinoid Nullstellensatz in the same theorem makes every residue field
$B/\mathfrak m$ finite over $K$; it is already complete.  Therefore every
point of $Z$ is rigid, not necessarily $K$-rational.  Put
\begin{equation}
 \label{eq:p21p3-componentwise-conductor-preimage}
                         Z_i:=\bar\nu_i^{-1}(Z).
\end{equation}
A finite morphism induces a finite extension of completed residue fields,
so every point of $Z_i$ is rigid.  Properness of a finite map and local
finiteness of $Z$ show that $Z_i$ is locally finite.

A finite algebra over a strict affinoid algebra is strict by
\cite[Theorem~4.4]{Tate71Rigid}; hence $\widetilde X_i$ is strict.  The maps
$X_i\to X_{\mathrm{red}}$ and
$\widetilde X_i\to X_i$ are finite and have empty relative boundary.
Applying \cite[(11.4.1.3)]{CLD25} twice shows that
$\widetilde X_i$ is boundaryless.  It is good and separated, normal by
construction, and connected because it normalizes one irreducible
component.

We shall apply the already verified normal rigid-exception proposition on
\begin{equation}
 \label{eq:p21p3-componentwise-normal-open}
                         U_i:=\widetilde X_i\setminus Z_i.
\end{equation}
For completeness, we check all of its topological hypotheses.  At a rigid
point of a normal analytic curve the normalization fibre has cardinality
one, so the point is unibranch by
\cite[Proposition~3.5.2]{Ducros24CurveStructure}.  Every connected analytic
curve is a graph and is connected by injective arcs by
\cite[Theorem~3.5.1 and \S1.3.5]{Ducros24CurveStructure}, whereas an
interior point of an injective arc is pluribranch by
\cite[\S1.3.12]{Ducros24CurveStructure}.  It follows that an injective arc
whose endpoints avoid $Z_i$ cannot meet $Z_i$; hence $U_i$ is connected.
It is nonempty as well: if one starts at a point of $Z_i$, its unique branch
contains an injective subarc with a pluribranch interior point, and that
point is not rigid and hence not in $Z_i$.

The space $U_i$ is therefore nonempty, connected, good, separated, strict,
normal, and boundaryless.  The set
\[
                         E_i':=E_i\setminus Z_i\subset U_i
\]
is locally finite and rigid, while
\[
 U_i\setminus E_i'
   =\widetilde X_i\setminus(E_i\cup Z_i)
\]
is an open subspace of the rig-smooth curve
$\widetilde X_i\setminus E_i$ and is therefore rig-smooth.  Thus every
hypothesis of \cref{prop:p21p3-rigid-exception-extension} is verified for
the pair $(U_i,E_i')$.

\paragraph{Generic multiplicities.}
By \cite[Definition~13.2.5 and \S13.5.2]{CLD25},
\[
                   [X]=\sum_i m_i[X_i],
 \qquad
                   [X_{\mathrm{red}}]=\sum_i[X_i].
\]
Integration agrees with the current attached to the fundamental cycle by
\cite[Proposition~16.6.3]{CLD25}.  The normalization is finite and has
calibrated degree one over the normal locus.  Its exceptional image lies in
$Z$, and every smooth $(1,1)$-test-form germ vanishes at a rigid point by
\cite[Proposition~3.3.6 and Proposition~15.2.1(a)]{CLD25}.  Consequently
the cycle identities hold as current identities on all of $X$:
\[
       \delta_X=\sum_i m_i(\nu_i)_*\delta_{\widetilde X_i},
 \qquad
       \iota_*\delta_{X_{\mathrm{red}}}
          =\sum_i(\nu_i)_*\delta_{\widetilde X_i}.
\]
This is \eqref{eq:p21p3-componentwise-fundamental-cycles}.  Local finiteness
of the components makes both sums locally finite.  If all $m_i$ equal one,
the two currents are equal.  Conversely, restrict an asserted equality to a
nonempty rig-smooth open subset of $X_i$ meeting neither $Z$ nor another
component.  The nonzero integration current there gives $m_i=1$.  Since
$m_i$ is the length of the generic local ring along $X_i$, this is
equivalent to generic reducedness along every component.  This proves
\cref{itm:p21p3-componentwise:integration}.

\paragraph{Current classification and uniqueness.}
Let $T$ be a positive $\mathrm d'$- and $\mathrm d''$-closed current of
bidegree $(0,0)$ on $X$.  By
\eqref{eq:p21p3-componentwise-current-isomorphism}, there is a unique
positive closed current $T_{\mathrm{red}}$ on $X_{\mathrm{red}}$ with
$T=\iota_*T_{\mathrm{red}}$.  The open set
$X_{\mathrm{red}}\setminus Z$ is the disjoint union of the normal component
loci $X_i\setminus Z$, and
\[
 \bar\nu_i|_{U_i}:U_i\xrightarrow{\ \sim\ }X_i\setminus Z.
\]
Transport the restriction of $T_{\mathrm{red}}$ on $X_i\setminus Z$ to a
current $T_i$ on $U_i$.  Restriction to an open subspace preserves
positivity and closedness by \cite[\S7.1.6 and \S\S7.7.1--7.7.3]{CLD25}.
The verified classification
\cref{itm:p21p3-rigid-exception:classification} applied to $(U_i,E_i')$
therefore gives a unique $c_i\in\mathbf R_{\geq0}$ such that
\begin{equation}
 \label{eq:p21p3-componentwise-classification-off-Z}
                            T_i=c_i\delta_{U_i}.
\end{equation}

Write
$\bar\delta_i:=(\bar\nu_i)_*\delta_{\widetilde X_i}$ on
$X_{\mathrm{red}}$, and set
\[
                         S:=T_{\mathrm{red}}-\sum_i c_i\bar\delta_i.
\]
The sum is locally finite, and
\eqref{eq:p21p3-componentwise-classification-off-Z} says that $S$ vanishes
on $X_{\mathrm{red}}\setminus Z$.  Let $\alpha$ be a compactly supported
smooth $(1,1)$-test form.  A rigid point $z$ has $d_K(z)=0$, since
$\mathscr H(z)/K$ is finite and both terms in the dimension formula
\cite[(11.2.3.1)]{CLD25} vanish.  The analytic support theorem
\cite[Proposition~15.2.1(a)]{CLD25} gives
\[
 \operatorname{Supp}(\alpha)
    \subseteq\{x:d_K(x)\geq1\}
    \subseteq X_{\mathrm{red}}\setminus Z.
\]
Thus $\langle S,\alpha\rangle=0$, and $S=0$.  Pushing forward by $\iota$
proves \eqref{eq:p21p3-componentwise-classification}.

For uniqueness, fix $i$ and restrict an equality of two such sums to a
nonempty rig-smooth open subset of $X_i\setminus Z$ that meets no other
component.  Such a subset exists because
$U_i\setminus E_i'$ is nonempty by
\cref{itm:p21p3-rigid-exception:connected}.  The integration current on it
is nonzero, so the two coefficients of $\delta_i$ agree.  This proves both
uniqueness and \cref{itm:p21p3-componentwise:classification}.  Notice that
the zero-atom input has not been used.

\paragraph{Componentwise splitting of the curvature measures.}
Let $W,U,V$ be as in
\cref{itm:p21p3-componentwise:bedford-taylor}, and regard $W$ also as an
open subset of $X_{\mathrm{red}}$.  Put
\[
             \widetilde W_i:=\bar\nu_i^{-1}(W),
\]
and denote the restrictions of $\nu_i$ and $\bar\nu_i$ by
$\nu_{i,W}^X$ and $\nu_{i,W}^{\mathrm{red}}$, respectively.  The good
separated spaces $\widetilde X_i$ are locally differentially
separated by \cite[\S15.2.2 and Corollary~15.2.4]{CLD25}.  Pullback preserves
continuous psh functions by
\cite[Proposition~8.1.5]{CLD25}.  Hence
\begin{equation}
 \label{eq:p21p3-componentwise-measures}
 \mu_i^U:=\mathrm d'\mathrm d''[\bar\nu_i^*U],
 \qquad
 \mu_i^V:=\mathrm d'\mathrm d''[\bar\nu_i^*V]
\end{equation}
are positive Radon measures by \cite[Lemma~8.5.5]{CLD25}.  Put
$\sigma_i:=\mu_i^U-\mu_i^V$.

For every continuous function $F$ on $W$, multiply the two current
identities in \eqref{eq:p21p3-componentwise-fundamental-cycles} by $F$.
The product is defined by \cite[Example~7.2.12]{CLD25}, and direct testing
against a compactly supported form gives
\begin{align}
 [F]_X
 &=\sum_i m_i(\nu_{i,W}^X)_*
       \bigl(((\nu_{i,W}^X)^*F)\delta_{\widetilde W_i}\bigr),
 \label{eq:p21p3-componentwise-function-current-X}\\
 [F]_{X_{\mathrm{red}}}
 &=\sum_i(\nu_{i,W}^{\mathrm{red}})_*
       \bigl(((\nu_{i,W}^{\mathrm{red}})^*F)
                              \delta_{\widetilde W_i}\bigr).
 \label{eq:p21p3-componentwise-function-current-red}
\end{align}
Finite proper pushforward commutes with the two current differentials by
transposition; for continuous psh functions this is also
\cite[Proposition~8.5.6]{CLD25}.  Differentiating
\eqref{eq:p21p3-componentwise-function-current-X} for $F=U,V$ shows that
the left equality in
\eqref{eq:p21p3-componentwise-curvature-equivalence} is equivalent to
\begin{equation}
 \label{eq:p21p3-componentwise-weighted-sum}
                 \sum_i m_i(\nu_{i,W}^X)_*\sigma_i=0.
\end{equation}
Since $\nu_{i,W}^X=\iota\circ\nu_{i,W}^{\mathrm{red}}$ and $\iota_*$ is
injective by \eqref{eq:p21p3-componentwise-current-isomorphism}, this is
equivalent to
\begin{equation}
 \label{eq:p21p3-componentwise-weighted-sum-red}
       \sum_i m_i(\nu_{i,W}^{\mathrm{red}})_*\sigma_i=0
       \qquad\text{on }X_{\mathrm{red}}.
\end{equation}
Likewise the right equality is equivalent to
\begin{equation}
 \label{eq:p21p3-componentwise-unweighted-sum}
                 \sum_i(\nu_{i,W}^{\mathrm{red}})_*\sigma_i=0.
\end{equation}

Set $Z_W:=Z\cap W$ and
$Z_{i,W}:=Z_i\cap\widetilde W_i$.  Away from $Z_W$, the component images in
$X_{\mathrm{red}}$ are disjoint and every $\nu_{i,W}^{\mathrm{red}}$ is an
analytic isomorphism onto its component image.  Restricting
\eqref{eq:p21p3-componentwise-weighted-sum-red} to the image of one component,
and using $m_i>0$, gives
\begin{equation}
 \label{eq:p21p3-componentwise-support}
                    \operatorname{Supp}(\sigma_i)\subseteq Z_{i,W}.
\end{equation}
Exactly the same conclusion follows from
\eqref{eq:p21p3-componentwise-unweighted-sum}.

Every point of $Z_{i,W}$ is rigid.  The space $\widetilde X_i$ is
boundaryless, pure one-dimensional, and normal, so every such point lies in
its interior and has a normal analytic local ring.  Apply
\cref{lem:p21p3-componentwise-normal-rigid-zero-atom} on the open curve
$\widetilde W_i$ to each of the two continuous psh functions in
\eqref{eq:p21p3-componentwise-measures}.  It gives
\[
                         \mu_i^U(\{z\})
                         =\mu_i^V(\{z\})=0
                         \qquad(z\in Z_{i,W}).
\]
We spell out why this kills the entire signed measure.  Since $Z_{i,W}$ is
locally finite, every $z\in Z_{i,W}$ has a neighborhood meeting that set
only at $z$.  On such a neighborhood,
\eqref{eq:p21p3-componentwise-support} makes $\sigma_i$ a multiple of
$\delta_z$; its coefficient is
$\sigma_i(\{z\})=0$.  The measure also vanishes off $Z_{i,W}$.  Hence
\begin{equation}
 \label{eq:p21p3-componentwise-component-equality}
                            \sigma_i=0
                            \qquad\text{for every }i.
\end{equation}
Starting from either the weighted or the unweighted sum therefore gives
\eqref{eq:p21p3-componentwise-component-equality}; conversely that equality
implies both sums.  This proves the exact equivalence
\eqref{eq:p21p3-componentwise-curvature-equivalence}.  In particular,
neither different generic multiplicities nor component intersections can
hide cancellation.

Finally put $h=U-V$ and use
\eqref{eq:p21p3-componentwise-classification}.  Multiplication by $h$ is
defined because $T$ is positive and $h$ is locally bounded.  On compact
test supports the component sum is finite.  Proper pushforward,
\eqref{eq:p21p3-componentwise-component-equality}, and linearity give
\begin{align*}
 \mathrm d'\mathrm d''(hT|_W)
 &=\sum_i c_i(\nu_{i,W}^X)_*
       \mathrm d'\mathrm d''[\bar\nu_i^*U-\bar\nu_i^*V]\\
 &=\sum_i c_i(\nu_{i,W}^X)_*(\mu_i^U-\mu_i^V)=0.
\end{align*}
This is \eqref{eq:p21p3-componentwise-bt-conclusion}.  The proof uses only
continuous psh pairs and makes no claim about the stronger discontinuous
quantifier in \cite[Definition~8.6.2]{CLD25}.
\end{proof}
